% 01_introduction.tex — Introduction and contributions

\section{Introduction}\label{sec:introduction}

\paragraph{Problem motivation}
Substantial funding flows into computer-science research each year, yet project post-mortems frequently cite months lost to peripheral engineering, late pivots, and analysis paralysis.
A lack of systematic guidance leaves even well-resourced teams mired in uncertainty.

\paragraph{Prevailing assumption}
The community tacitly assumes that close proximity to seasoned mentors suffices for transmitting the craft of discovery.
Apprenticeship produces excellence, but it does not scale: heuristics remain siloed and training outcomes vary widely.

\paragraph{Proposed framework}
We formalise an alternative: a decision loop that (i) orders open questions by impact and uncertainty and (ii) resolves the top-ranked question via the fastest informative experiment.
This loop replaces folkloric judgment with an explicit protocol.

\paragraph{Implementation}
To operationalise the loop, we release an open-source toolbox containing assumption maps, sprint planners, and a velocity diagnostic that tracks learning per unit time.
These artefacts translate abstract principles into daily practice.

\paragraph{Evaluation}
We applied the toolbox to a graduate machine-learning project, an industry research prototype, and a seed-stage startup.
Across all settings, the loop cut time-to-insight by a factor of two to four relative to retrospective baselines.

\paragraph{Implications}
Integrating the loop into curricula, lab meetings, and corporate R\&D programmes promises faster iteration, clearer contribution statements, and more predictable research outcomes.

% ------------------------------------------------
\subsection*{Contributions}
\begin{itemize}
  \item Consolidate Bernstein's lecture guidance into a formal, citable framework for high-uncertainty research.
  \item Provide implementation artefacts---assumption maps, sprint planners, velocity diagnostics---for immediate adoption.
  \item Demonstrate two-to-four-fold faster insight cycles across academic and industrial case studies.
  \item Offer curricular and organisational recommendations for scalable, systematic research training.
\end{itemize}

\subsection*{Roadmap}
Section~\ref{sec:abstract_blueprint} presents the six-move abstract blueprint.
Section~\ref{sec:ml_research_advice} bridges production hygiene and frontier-scale ML research.
Section~\ref{sec:elon_five_steps} describes Musk's five-step protocol.
Section~\ref{sec:cs197_insights} formalises the Velocity--Vectoring framework from CS197.
Section~\ref{sec:relatedwork} situates our approach within project-management and design-thinking literature. \todo{Create sec:relatedwork section}
Section~\ref{sec:discussion} discusses limitations and future work. \todo{Create sec:discussion section}
Section~\ref{sec:conclusion} concludes. \todo{Create sec:conclusion section}
